\chapter{Task 3 -- Aging Interpolation}
\label{chap:task3}

\paragraph{Riassunto del capitolo.} Il terzo e più difficile task
chiede: caratterizzata la cella a quattro livelli di invecchiamento,
possiamo prevedere la caratterizzazione completa del quinto -- mai
osservato -- usando solo gli altri quattro? Lo splitter \acrshort{logo}
sull'\emph{Aging} simula esattamente questo. KNN distance-weighted
vince con $R^{2} = 0.986$ perché \emph{interpola} naturalmente nello
spazio scalato delle feature, mentre Random Forest e Gradient Boosting
sono limitati dalla loro incapacità di estrapolare oltre l'inviluppo dei
target visti. La performance degrada al bordo termico
(\SI{47.5}{\celsius}) ma resta sopra $0.95$.

\section{Problem statement}

\textbf{Domanda di business.} Se in laboratorio caratterizziamo le celle a
\emph{quattro} stadi di invecchiamento, possiamo prevedere la
caratterizzazione \emph{completa} del quinto stadio mai osservato?

\textbf{Domanda formale.} Si esclude un intero livello $Aging^*$ dal
training. Sull'addestramento contiamo $4 \times 8 \times 5 \times 49
\approx 7845$ misure; sul test, l'aging escluso conta
$1 \times 8 \times 5 \times 49 \approx 1960$ misure -- 40 curve di Nyquist
da ricostruire (8 temperature $\times$ 5 SOC).

\textbf{Difficoltà.} Maggiore di Task 1: lì il modello vedeva durante
l'addestramento \emph{altre} curve allo stesso $Aging^*$ (per le altre
temperature) e altre $Temperature^*$ (per gli altri aging). Qui $Aging^*$
non viene mai osservato.

\section{Setup sperimentale}

\begin{itemize}
  \item \textbf{Aging escluso (default):} 2 (centrale, per testare
        l'interpolazione anziché l'estrapolazione).
  \item \textbf{Feature:} le stesse nove di \cref{tab:regression-features}.
  \item \textbf{Target:} $(\mathrm{Re}(Z), \mathrm{Im}(Z))$, multi-output.
  \item \textbf{Split:} maschera booleana sull'attributo $Aging$.
  \item \textbf{Funzione di produzione:}
        \codefile{src/models/task3\_aging.py::run\_aging\_interpolation}
        $\to$ \texttt{AgingInterpResult} con metriche per-temperatura.
\end{itemize}

\section{Modelli candidati}

Task 3 utilizza esattamente la stessa funzione factory di Task 1 \\
(\codefile{src/models/registry.py::regression\_models}), perché lo
spazio delle feature e il target sono identici. Gli iperparametri di
default sono quelli della~\cref{tab:task1-models}: Ridge $\alpha=1$,
Random Forest $n_\text{estim}=200$ con $d_\text{max}=15$ e
$\text{leaves} \geq 2$, Gradient Boosting $n_\text{estim}=150$, $d=5$,
$\eta=0.1$ in \texttt{MultiOutputRegressor}, KNN $k=10$ con weight
\texttt{distance}, e Bagging (Ridge) con 30 estimator e sub-sample
$0.8$.

\section{Risultati}

\begin{table}[H]
  \centering
  \caption{Benchmark Task 3 (esclusione: $Aging = 2$). Anche qui il
  baseline \emph{mean predictor} è il punto di riferimento $R^{2}\!\approx\!0$
  contro cui leggere i guadagni.}
  \label{tab:task3-benchmark}
  \begin{tabular}{lrrr}
    \toprule
    \textbf{Modello} & $R^{2}$ & \acrshort{mse} & \acrshort{mae} \\
    \midrule
    K-Nearest Neighbors & \textbf{0.9862} & \textbf{0.0049} & \textbf{0.0449} \\
    Gradient Boosting   & 0.9795 & 0.0227 & 0.0932 \\
    Random Forest       & 0.9753 & 0.0287 & 0.0950 \\
    Ridge               & 0.7293 & 0.1919 & 0.2930 \\
    Bagging Ridge       & 0.7289 & 0.1924 & 0.2936 \\
    \midrule
    \emph{Mean predictor (baseline)} & $\approx 0.0$ & $\mathrm{Var}(y_\text{test})$ & $\mathrm{MAD}(y_\text{test})$ \\
    \bottomrule
  \end{tabular}
\end{table}

\textbf{KNN vince.} Il motivo è strutturale: nello spazio standardizzato
delle nove feature, predire $Aging = 2$ significa cercare i vicini più
prossimi nello \emph{strato} di feature
$(Aging \in \{0, 1, 3, 4\}, \dots)$, e KNN con weight \texttt{distance}
restituisce una media pesata che è effettivamente un'\emph{interpolazione
locale}. Random Forest e Gradient Boosting, al contrario, predicono solo
medie di foglie osservate durante il training: \emph{non} possono
estrapolare a un valore di $Aging$ mai visto, e le loro predizioni si
appiattiscono verso il centro convesso dei target visti.

\subsection{Analisi per temperatura}

\codefile{run\_aging\_interpolation} restituisce un DataFrame
\texttt{per\_temperature} con $R^{2}$ e \acrshort{mae} per ciascun
livello di temperatura. I numeri di seguito si riferiscono al modello
KNN sul fold di default.

\begin{table}[H]
  \centering
  \caption{Performance per temperatura del KNN su Task 3.}
  \label{tab:task3-per-temp}
  \begin{tabular}{rrr}
    \toprule
    Temperature [\si{\celsius}] & $R^{2}$ & \acrshort{mae} [\milliohm] \\
    \midrule
    20.0 & 0.9894 & 0.0643 \\
    22.5 & 0.9842 & 0.0714 \\
    25.0 & 0.9860 & 0.0439 \\
    27.5 & 0.9872 & 0.0332 \\
    30.0 & 0.9815 & 0.0489 \\
    35.0 & 0.9787 & 0.0254 \\
    40.0 & 0.9739 & 0.0300 \\
    47.5 & \textbf{0.9566} & 0.0422 \\
    \bottomrule
  \end{tabular}
\end{table}

L'$R^{2}$ è massimo nelle temperature centrali (\SI{27.5}{\celsius}:
$0.9872$) e degrada al bordo superiore (\SI{47.5}{\celsius}:
$0.9566$): è un effetto di \emph{boundary} atteso, oltre il bordo non
ci sono campioni che possano sostenere l'interpolazione e KNN ricade su
vicini più ``lontani'' nello spazio scalato. La \acrshort{mae} resta
sempre sotto $\SI{0.072}{\milliohm}$ in modulo.

\section{Hyperparameter tuning}

\textbf{Splitter:} \acrshort{logo} sul gruppo $Aging$ (5 fold, ognuno
tiene fuori un livello di aging — niente sottocampionamento perché la
dimensione è già piccola).

\textbf{Grid (da \codefile{config.yaml}):}
\begin{lstlisting}[language=YAML]
knn:
  n_neighbors: [5, 10, 15]
  weights:     [uniform, distance]
\end{lstlisting}

Sei combinazioni $\times$ 5 fold = 30 fit. La metrica di selezione è
la media \acrshort{mse} sui 5 fold \acrshort{logo}. Rispetto al
default registry $\{k=10, \text{weights}=\text{distance}\}$ il tuning
identifica configurazioni con $k$ leggermente più piccolo, ma il
guadagno sul fold di test non è statisticamente distinguibile. Vale
lo stesso caveat statistico: con 5 fold informativamente molto diversi
e una sola valutazione finale, senza intervallo di confidenza
calcolato per bootstrap (cfr.~\cref{sec:bootstrap}) non possiamo
concludere che il tuning peggiori. Manteniamo la baseline.

\newpage
\section{Visualizzazioni}

\begin{figure}[htbp]
  \centering
  \includegraphics[width=0.9\linewidth]{task3_metrics_bar.pdf}
  \caption{$R^{2}$ sul fold di test per i cinque modelli candidati.
  KNN distanza-pesato è il solo a interpolare correttamente; i modelli
  lineari saturano sotto $0.73$.}
  \label{fig:task3-metrics-bar}
\end{figure}

\begin{figure}[htbp]
  \centering
  \includegraphics[width=0.90\linewidth]{task3_per_temp.pdf}
  \caption{$R^{2}$ e MAE per livello di temperatura del modello KNN
  vincitore. La performance degrada al bordo superiore della finestra
  termica (\SI{47.5}{\celsius}, $R^{2} = 0.957$) per effetto di
  \emph{boundary}.}
  \label{fig:task3-per-temp}
\end{figure}

\begin{figure}[htbp]
  \centering
  \includegraphics[width=0.68\linewidth]{task3_grid_sample.pdf}
  \caption{Quattro curve di Nyquist riconcostruite (rosso) sovrapposte
  alle vere (nero), prese a due temperature e due SOC diversi. Pur non
  avendo mai osservato $Aging = 2$ in addestramento, il modello
  ricostruisce correttamente sia il semicerchio sia la coda diffusiva.}
  \label{fig:task3-sample}
\end{figure}

\begin{figure}[htbp]
  \centering
  \includegraphics[width=0.78\linewidth]{task3_error_distribution.pdf}
  \caption{Istogramma degli errori assoluti
  $\sqrt{\Delta\mathrm{Re}^2 + \Delta\mathrm{Im}^2}$ sul fold di test.
  Il 95\textsuperscript{o} percentile cade sotto i \SI{0.2}{\milliohm},
  ben dentro la dispersione fisiologica delle misure.}
  \label{fig:task3-error}
\end{figure}
